\chapter{Proces Wienera: własności probabilistyczne i geometria trajektorii}
\label{ch:wiener}
\paragraph{Odsyłacze międzytomowe.} Rachunek całkowy względem Wienera zaczyna się w tomie III, rozdziale~\ref{III-ch:wiener-integral}; zastosowania symulacyjne znajdują się w tomie IV, rozdziale~\ref{IV-ch:monte-carlo}.


\section{Od sum niezależnych przyrostów do ruchu Browna}
Rozważmy błąd pomiaru, który narasta jako suma bardzo wielu drobnych,
w przybliżeniu niezależnych zakłóceń. W granicy czasowej spodziewamy
się procesu o przyrostach gaussowskich, którego wariancja rośnie
proporcjonalnie do długości przedziału. Tak powstaje matematyczna
motywacja standardowego procesu Wienera, czyli ruchu Browna.
Dla każdego $t$ mamy rozkład normalny, ale znacznie silniejszy
jest warunek \emph{niezależności przyrostów na rozłącznych
przedziałach}. To właśnie on odróżnia Wienera od dowolnego
procesu gaussowskiego.

\begin{definicja}[Standardowy proces Wienera]
Niech $(\Omega,\calF,\Prob)$ będzie przestrzenią probabilistyczną.
Proces rzeczywisty $W=(W_t)_{t\ge0}$ nazywamy
\emph{standardowym procesem Wienera} (\emph{standardowym ruchem Browna}),
jeżeli:
\begin{enumerate}
\item $W_0=0$ prawie na pewno;
\item jego trajektorie $t\mapsto W_t(\omega)$ są ciągłe dla prawie
wszystkich $\omega$;
\item dla dowolnych $0\le t_0<t_1<\cdots<t_n$ przyrosty
$W_{t_1}-W_{t_0},\ldots,W_{t_n}-W_{t_{n-1}}$
są wzajemnie niezależne;
\item $W_t-W_s\sim\Normal(0,t-s)$ dla wszystkich $0\le s<t$.
\end{enumerate}
Jeżeli ponadto zadano filtrację $(\calF_t)_{t\ge0}$, to $W$
nazywamy procesem Wienera \emph{względem tej filtracji},
gdy jest do niej adaptowany oraz $W_t-W_s$ jest niezależny
od $\calF_s$ dla wszystkich $s<t$.
\end{definicja}

Rozdział~\ref{ch:konstrukcja} dowiódł już istnienia takiego
procesu na pewnej przestrzeni probabilistycznej.
W dalszym ciągu możemy bez straty ogólności przyjąć,
że wszystkie jego trajektorie są ciągłe: modyfikujemy proces
na jednym zbiorze zerowym, przypisując na nim trajektorię zerową.
Takie ustalenie nie zmienia rozkładów skończeniewymiarowych.

\begin{twierdzenie}[Równoważna charakteryzacja gaussowska]
Niech $X=(X_t)_{t\ge0}$ będzie rzeczywistym procesem
o prawie na pewno ciągłych trajektoriach. Wówczas $X$
jest standardowym procesem Wienera wtedy i tylko wtedy,
gdy jest procesem gaussowskim i dla wszystkich $s,t\ge0$
zachodzi
\[
\E X_t=0,\qquad \Cov(X_s,X_t)=\min(s,t).
\]
\end{twierdzenie}
\begin{proof}
Jeżeli $X=W$ jest procesem Wienera, to dla
$0=t_0<t_1<\cdots<t_n$ wektor $(W_{t_1},\ldots,W_{t_n})$
jest liniowym obrazem wektora niezależnych przyrostów normalnych,
więc jest gaussowski. Przy $s\le t$ korzystamy z rozkładu
$W_t=W_s+(W_t-W_s)$ oraz niezależności przyrostu od $W_s$:
$\Cov(W_s,W_t)=\Var(W_s)=s$. Średnie wynoszą zero.

Odwrotnie, $\Var(X_0)=K(0,0)=0$ i $\E X_0=0$,
a więc $X_0=0$ p.n. Dla $s<t$
$\Var(X_t-X_s)=t+s-2s=t-s$ i średnia przyrostu
jest zero. Jeżeli $[a,b]$ i $[c,d]$ są rozłączne oraz $b\le c$,
to $\Cov(X_b-X_a,X_d-X_c)=0$, co wynika z podstawienia
czterech wartości funkcji $\min(\cdot,\cdot)$.
Dowolne skończone rodziny takich przyrostów są wspólnie
gaussowskie, zatem ich zerowe kowariancje oznaczają niezależność.
Dostajemy wszystkie punkty definicji.
\end{proof}

\begin{przyklad}[Rozkład kilku obserwacji]
Dla $0<s<t$ wektor $(W_s,W_t)$ jest normalny dwuwymiarowo
ze średnią zero i macierzą kowariancji
$\left(\begin{smallmatrix}s&s\\s&t\end{smallmatrix}\right)$.
Jego współrzędne są zależne: $\Cov(W_s,W_t)=s$.
Jednak $W_s$ i $W_t-W_s$ są niezależne.
Dla $0<t_1<t_2<t_3$ mamy na przykład
\[
\Prob(W_{t_3}-W_{t_2}>a\mid W_{t_1}=x)
=1-\Phi\left(\frac{a}{\sqrt{t_3-t_2}}\right),
\]
gdzie równość rozumiana jest jako odpowiednia wersja
regularnego prawdopodobieństwa warunkowego.
Warunek na $W_{t_1}$ nie wpływa na późniejszy rozłączny przyrost.
Nie należy tego mylić z pytaniem o rozkład $W_{t_3}$,
który oczywiście zależy od stanu w chwili $t_1$.
\end{przyklad}

\section{Skalowanie, symetria i niezależność przyszłości}
\begin{twierdzenie}[Symetrie i samopodobieństwo]
Niech $W$ będzie procesem Wienera, a $c>0$.
Procesy $(c^{-1/2}W_{ct})_{t\ge0}$ i
$(-W_t)_{t\ge0}$ mają to samo prawo na przestrzeni
ciągłych trajektorii co $W$. Dla każdego $T>0$ proces
$(W_T-W_{T-t})_{0\le t\le T}$ ma prawo procesu Wienera
ograniczonego do $[0,T]$.
\end{twierdzenie}
\begin{proof}
Każdy z wymienionych procesów ma ciągłe trajektorie,
wartość początkową zero i jest gaussowski.
Dla pierwszego procesu kowariancja wynosi
$c^{-1}\min(cs,ct)=\min(s,t)$; dla drugiego
jest taka sama. Dla trzeciego, gdy $0\le s\le t\le T$,
przyrosty $W_T-W_{T-s}$ oraz $W_T-W_{T-t}$ mają
kowariancję $s$: wspólnym składnikiem są przyrosty
na przedziale $[T-s,T]$. Z poprzedniego twierdzenia
i jednoznaczności prawa procesu otrzymujemy tezy.
\end{proof}

\begin{uwaga}[Jak rozumieć skalowanie?]
Jeśli skala czasu rośnie $c$ razy, typowe wychylenie
rośnie $\sqrt c$ razy, a nie $c$ razy.
To odpowiada zależności $\Var(W_t)=t$.
Samopodobieństwo stanowi podstawę późniejszej zasady
niezmienniczości Donskera, wiążącej spacery losowe z
ruchem Browna.
\end{uwaga}

\begin{twierdzenie}[Własność Markowa i półgrupa przejścia]
Niech $W$ będzie procesem Wienera względem filtracji
$(\calF_t)_{t\ge0}$. Dla dowolnych $s,t\ge0$
i ograniczonej mierzalnej funkcji $f:\R\to\R$ zachodzi
\[
\E[f(W_{s+t})\mid\calF_s]
=(P_tf)(W_s),\qquad
(P_tf)(x):=\int_\R f(x+y)p_t(y)\,\dd\lambda_1(y),
\]
gdzie $p_t(y)=(2\pi t)^{-1/2}\exp(-y^2/(2t))$ dla $t>0$,
a $P_0f=f$. Ponadto $P_{s+t}=P_sP_t$.
\end{twierdzenie}
\begin{proof}
Zapiszmy $W_{s+t}=W_s+Z$, gdzie $Z=W_{s+t}-W_s$
jest niezależna od $\calF_s$ i ma rozkład $\Normal(0,t)$.
Dla funkcji $f(x)=\ind{\{x\in A\}}$ własność
warunkowania niezależnej zmiennej daje żądany wzór.
Przez liniowość uzyskujemy go dla funkcji prostych, a
następnie przez zbieżność monotoniczną i rozkład na
części dodatnią i ujemną dla wszystkich ograniczonych
funkcji mierzalnych.
Aby udowodnić własność półgrupy, bierzemy niezależne
$Z_s\sim\Normal(0,s)$ i $Z_t\sim\Normal(0,t)$.
Ich suma ma rozkład $\Normal(0,s+t)$. Stąd
$P_s(P_tf)(x)=\E f(x+Z_s+Z_t)=P_{s+t}f(x)$;
przejścia $s=0$ lub $t=0$ są oczywiste.
\end{proof}

\section{Mocna własność Markowa: można zacząć od losowej chwili}
Dla deterministycznej chwili $s$ przyszły przyrost
$W_{s+t}-W_s$ jest niezależny od $\calF_s$.
Prawdziwie nowa treść pojawia się przy \emph{losowej}
chwili $\tau$, np. pierwszym trafieniu w poziom $a$.
Nie można po prostu podstawić $s=\tau(\omega)$ w definicji
niezależności: musimy najpierw przybliżyć czas losowy
czasami przyjmującymi przeliczalnie wiele wartości.

\begin{twierdzenie}[Mocna własność Markowa procesu Wienera]
Niech $W$ będzie procesem Wienera względem filtracji
$(\calF_t)_{t\ge0}$, a $\tau$ niech będzie względem niej
czasem zatrzymania takim, że $\Prob(\tau<\infty)=1$.
Wtedy proces $(W_{\tau+t}-W_\tau)_{t\ge0}$ jest
standardowym procesem Wienera i jest niezależny od
sigma-ciała $\calF_\tau$.
\end{twierdzenie}
\begin{proof}
Dla $n\ge1$ niech $h_n=2^{-n}$ i
$\tau_n=h_n\lceil\tau/h_n\rceil$.
Ponieważ $\{\tau_n\le t\}$ jest zdarzeniem określonym
przez warunek $\{\tau\le h_n\lfloor t/h_n\rfloor\}$,
czas $\tau_n$ jest czasem zatrzymania, przyjmuje wartości
w siatce $h_n\N_0$, spełnia $\tau\le\tau_n<\tau+h_n$
i zbiega do $\tau$ prawie na pewno.

Niech $A\in\calF_\tau$. Z nierówności $\tau\le\tau_n$
wynika $A\in\calF_{\tau_n}$. Ustalmy
$0\le t_1<\cdots<t_m$ i $u=(u_1,\ldots,u_m)\in\R^m$.
Dla zdarzenia $A\cap\{\tau_n=kh_n\}\in\calF_{kh_n}$
niezależność przyrostów po chwili $kh_n$ od tej
sigma-algebry daje
\begin{align*}
&\E\left[\ind{A\cap\{\tau_n=kh_n\}}
\exp\left(i\sum_{j=1}^m u_j(W_{\tau_n+t_j}-W_{\tau_n})\right)\right]\\
&\qquad=\Prob(A\cap\{\tau_n=kh_n\})
\exp\left(-\frac12\sum_{j,l=1}^m u_ju_l\min(t_j,t_l)\right).
\end{align*}
Sumujemy po $k\ge0$; jest to dozwolone przez
zbieżność zdominowaną, gdyż moduł każdego czynnika
wykładniczego nie przekracza $1$. Otrzymujemy
\[
\E\left[\ind A\exp\left(i\sum_j u_j
(W_{\tau_n+t_j}-W_{\tau_n})\right)\right]
=\Prob(A)\exp\left(-\frac12\sum_{j,l}
 u_ju_l\min(t_j,t_l)\right).
\]
Ciągłość trajektorii i $\tau_n\to\tau$ zapewniają
$W_{\tau_n+t_j}-W_{\tau_n}\to W_{\tau+t_j}-W_\tau$
prawie na pewno. Ponownie stosujemy zbieżność zdominowaną.
Otrzymana równość dla wszystkich $u$ oznacza, z jednoznaczności
funkcji charakterystycznych, że wektor przyszłych przyrostów
ma standardowy rozkład gaussowski i jest niezależny
od każdego $A\in\calF_\tau$. Niezależność
całej ciągłej trajektorii od $\calF_\tau$ wynika z faktu,
że borelowskie sigma-ciało przestrzeni ciągłych trajektorii
jest generowane przez odczyty w chwilach wymiernych.
Ciągłość przyszłych trajektorii wynika z ciągłości $W$.
\end{proof}

\begin{przyklad}[Co znaczy ``zapominanie'' w losowej chwili?]
Niech $a>0$ i $\tau_a=\inf\{t\ge0:W_t=a\}$.
Na zdarzeniu $\{\tau_a<\infty\}$ obserwator zna
całą przeszłą drogę oraz fakt, że cząstka właśnie dotarła
do $a$. Mocna własność Markowa mówi, że przyszły \emph{przyrost}
$W_{\tau_a+t}-a$ jest niezależnym, nowym ruchem Browna.
Nie oznacza to niezależności samego położenia $W_{\tau_a+t}$
od $W_{\tau_a}$, bo położenie zawiera znany punkt startowy.
\end{przyklad}

\section{Regularność trajektorii: ciągłość bez gładkości}
Dla ustalonych $s<t$ przyrost ma skalę $\sqrt{t-s}$,
a nie $t-s$. Ta pozornie drobna różnica ma zasadniczy
skutek: trajektorie są ciągłe, ale prawie na pewno nie
mają pochodnej w żadnym punkcie.

\begin{twierdzenie}[Regularność H\"oldera poniżej $1/2$]
Dla każdego $T>0$ i standardowego procesu Wienera $W$
zachodzi z prawdopodobieństwem $1$: dla każdego
$\gamma\in(0,1/2)$ istnieje losowa stała
$C_{\gamma,T}(\omega)<\infty$ taka, że
\[
|W_t(\omega)-W_s(\omega)|
\le C_{\gamma,T}(\omega)|t-s|^\gamma,
\qquad s,t\in[0,T].
\]
\end{twierdzenie}
\begin{proof}
Dla każdego $p>2$ i $s,t\in[0,T]$ mamy
$\E|W_t-W_s|^p=c_p|t-s|^{p/2}$.
Kryterium ciągłości Kołmogorowa daje modyfikację
o trajektoriach H\"oldera każdego rzędu
$\gamma<1/2-1/p$. Ponieważ zarówno $W$, jak i ta
modyfikacja mają ciągłe trajektorie, są nierozróżnialne;
możemy więc przypisać tę własność $W$.
Wybierając przeliczalny ciąg $p_k\to\infty$,
przecinamy zbiory o prawdopodobieństwie $1$ i
otrzymujemy jednocześnie wszystkie $\gamma<1/2$.
\end{proof}

\section{Wariacja kwadratowa: niezwykła wielkość trajektorii}
Dla gładkiej funkcji $f$ i podziałów o coraz mniejszej
średnicy suma kwadratów przyrostów $\sum(\Delta f)^2$
dąży do zera. Dla Wienera dzieje się coś innego:
suma ta dąży do długości przedziału. To najważniejszy
mechanizm odróżniający rachunek Itô od klasycznego.

\begin{twierdzenie}[Wariacja kwadratowa procesu Wienera]
Niech $W$ będzie standardowym procesem Wienera, $T>0$, zaś
$\pi_n=\{0=t_{0,n}<\cdots<t_{m_n,n}=T\}$
ciągiem deterministycznych podziałów $[0,T]$ o średnicy
$|\pi_n|:=\max_k(t_{k+1,n}-t_{k,n})\to0$.
Wówczas
\[
Q(\pi_n):=\sum_{k=0}^{m_n-1}
(W_{t_{k+1,n}}-W_{t_{k,n}})^2
\longrightarrow T\quad\text{w }L^2.
\]
Jeżeli $\pi_n$ jest podziałem dyadycznym, to zbieżność
zachodzi również prawie na pewno.
\end{twierdzenie}
\begin{proof}
Oznaczmy $\Delta_{k,n}=W_{t_{k+1,n}}-W_{t_{k,n}}$
oraz $h_{k,n}=t_{k+1,n}-t_{k,n}$.
Przyrosty są niezależne, $\Delta_{k,n}\sim\Normal(0,h_{k,n})$,
a $\E\Delta_{k,n}^2=h_{k,n}$ i
$\E\Delta_{k,n}^4=3h_{k,n}^2$. Zatem
$\Var(\Delta_{k,n}^2)=2h_{k,n}^2$.
Dzięki niezależności
\[
\E Q(\pi_n)=\sum_k h_{k,n}=T,
\qquad
\E\bigl|Q(\pi_n)-T\bigr|^2
=2\sum_k h_{k,n}^2\le 2T|\pi_n|\to0.
\]
Dla siatek dyadycznych $|\pi_n|=T2^{-n}$.
Nierówność Czebyszewa daje, dla każdego $\varepsilon>0$,
$\Prob(|Q(\pi_n)-T|>\varepsilon)
\le 2T^2 2^{-n}/\varepsilon^2$.
Suma po $n$ jest skończona, więc z pierwszego lematu
Borela--Cantellego odchylenie większe niż $\varepsilon$
występuje tylko skończenie wiele razy prawie na pewno.
Stosując to do $\varepsilon=1,1/2,1/3,\ldots$
i przecinając przeliczalnie wiele zdarzeń pełnej miary,
uzyskujemy $Q(\pi_n)\to T$ p.n.
\end{proof}

\begin{wniosek}[Jednostajna zbieżność wariacji kwadratowej]
Dla podziałów dyadycznych $t_{k,n}=kT2^{-n}$
zdefiniujmy dla $t\in[0,T]$
\[
Q_n(t)=\sum_{k=0}^{2^n-1}
\bigl(W_{t\wedge t_{k+1,n}}-W_{t\wedge t_{k,n}}\bigr)^2.
\]
Wówczas $\sup_{0\le t\le T}|Q_n(t)-t|\to0$
prawie na pewno.
\end{wniosek}
\begin{proof}
Niech $h=T2^{-n}$ i
$Y_{k,n}=(W_{(k+1)h}-W_{kh})^2-h$.
Zmienne $Y_{k,n}$ są niezależne, wycentrowane i
$\sum_k\Var Y_{k,n}=2Th$.
Nierówność maksymalna Kołmogorowa daje
\[
\Prob\left(\max_{0\le j\le 2^n}
\left|\sum_{k<j}Y_{k,n}\right|>\varepsilon\right)
\le \frac{2Th}{\varepsilon^2}.
\]
Prawa strona jest sumowalna po $n$, więc maksimum
sum cząstkowych dąży do zera p.n. Gdy
$t\in[jh,(j+1)h]$, różnica między $Q_n(t)-t$
a sumą $\sum_{k<j}Y_{k,n}$ jest równa
$(W_t-W_{jh})^2-(t-jh)$. Jej moduł nie przekracza
$\max_j\sup_{u\in[jh,(j+1)h]}|W_u-W_{jh}|^2+h$.
Ta wielkość dąży do zera p.n. z jednostajnej
ciągłości trajektorii $W$ na $[0,T]$.
Połączenie obu oszacowań dowodzi tezy.
\end{proof}

\begin{wniosek}[Trajektorie nie mają skończonej wariacji całkowitej]
Dla każdego $T>0$ trajektorie $t\mapsto W_t(\omega)$
mają z prawdopodobieństwem $1$ nieskończoną wariację
całkowitą na $[0,T]$. Prawie na pewno jest tak nawet
na każdym przedziale o dodatniej długości.
\end{wniosek}
\begin{proof}
Gdyby pewna ciągła trajektoria $w$ miała skończoną
wariację całkowitą $V$, to dla dowolnego podziału
$\sum_k(w(t_{k+1})-w(t_k))^2
\le\max_k|w(t_{k+1})-w(t_k)|\,V$.
Jednostajna ciągłość $w$ zapewnia zbieżność maksimum
przyrostów do zera wraz ze średnicą podziału.
Stąd suma kwadratów dążyłaby do zera, co przeczy
poprzedniemu twierdzeniu, dającemu granicę $T>0$.
Aby uzyskać twierdzenie dla wszystkich nietrywialnych
przedziałów jednocześnie, stosujemy ten sam dowód do
przedziałów o wymiernych końcach i przecinamy
przeliczalnie wiele zbiorów pełnej miary.
Każdy przedział o dodatniej długości zawiera taki
wymierny podprzedział; wariacja na większym przedziale
nie może być mniejsza od wariacji na mniejszym.
\end{proof}

\begin{przyklad}[Dlaczego lewy i środkowy punkt przedziału dają inne całki?]
Niech $t_k=kT/n$. Tożsamość algebraiczna
$W_{t_{k+1}}^2-W_{t_k}^2
=2W_{t_k}\Delta_kW+(\Delta_kW)^2$
daje
\[
\sum_{k=0}^{n-1}W_{t_k}(W_{t_{k+1}}-W_{t_k})
=\frac12W_T^2-\frac12\sum_{k=0}^{n-1}(\Delta_kW)^2
\longrightarrow \frac12(W_T^2-T)
\quad\text{w }L^2.
\]
Z kolei suma z uśrednieniem wartości na końcach
$\frac12(W_{t_k}+W_{t_{k+1}})$ spełnia
\[
\sum_{k=0}^{n-1}\frac{W_{t_k}+W_{t_{k+1}}}{2}
 (W_{t_{k+1}}-W_{t_k})=\frac12W_T^2
\]
\emph{dokładnie} dla każdego $n$.
To algebraiczny początek rozróżnienia całki Itô i
Stratonovicha. Pełna konstrukcja obu całek pojawi
się w części o rachunku stochastycznym.
\end{przyklad}

\section{Dlaczego ruch Browna nie ma pochodnej?}
Brak skończonej wariacji całkowitej na każdym przedziale nie
wyklucza jeszcze pochodnej w pojedynczym punkcie. Potrzebny
jest silniejszy argument. Pokażemy, że na drobnej siatce
prawdopodobieństwo wystąpienia choćby trzech \emph{jednocześnie
bardzo małych} kolejnych przyrostów maleje wystarczająco szybko.
Gdyby trajektoria miała pochodną, takie trzy małe przyrosty
musiałyby się regularnie pojawiać blisko punktu różniczkowalności.

\begin{twierdzenie}[Nieróżniczkowalność trajektorii Wienera]
Niech $W$ będzie standardowym procesem Wienera i $T>0$.
Wówczas z prawdopodobieństwem $1$ trajektoria
$t\mapsto W_t(\omega)$ nie jest różniczkowalna
w żadnym punkcie $t\in(0,T)$ i nie ma skończonej
pochodnej jednostronnej w punktach $0,T$.
\end{twierdzenie}
\begin{proof}
Ustalmy liczbę całkowitą $M\ge1$ i siatkę dyadyczną
o kroku $h_n=T2^{-n}$. Niech $A_n(M)$ oznacza zdarzenie,
że istnieje $k\in\{0,\ldots,2^n-3\}$, dla którego
trzy kolejne przyrosty
$\Delta_{k,n}W,\Delta_{k+1,n}W,\Delta_{k+2,n}W$
mają moduł nie większy niż $8Mh_n$.
Każdy z tych przyrostów ma rozkład $\Normal(0,h_n)$,
a są one niezależne. Z oszacowania gęstości standardowej
normalnej przez $1/\sqrt{2\pi}$ dostajemy
\[
\Prob(|\Delta_{k,n}W|\le8Mh_n)
=\Prob(|G|\le8M\sqrt{h_n})
\le \frac{16M}{\sqrt{2\pi}}\sqrt{h_n}.
\]
Stąd, z niezależności trzech przyrostów i sumowania
prawdopodobieństw po najwyżej $T/h_n$ pozycjach,
\[
\Prob(A_n(M))\le C\,M^3T\sqrt{h_n},
\]
gdzie $C$ jest stałą liczbową. Szereg po $n$ jest zbieżny,
więc lemat Borela--Cantellego daje
$\Prob(A_n(M)\text{ nieskończenie często})=0$.

Przypuśćmy teraz, że ciągła trajektoria $w$ ma skończoną
pochodną w pewnym $t\in(0,T)$. Z definicji pochodnej
istnieją $M\in\N$ i $\delta>0$ takie, że
$|w(s)-w(t)|\le M|s-t|$ dla $|s-t|<\delta$.
Dla dostatecznie dużego $n$ istnieje $k$ z
$kh_n\le t<(k+1)h_n$, a punkty
$kh_n,(k+1)h_n,(k+2)h_n,(k+3)h_n$
leżą w $(t-\delta,t+\delta)\cap[0,T]$.
Dla $j=0,1,2$ nierówność trójkąta daje
\[
|w((k+j+1)h_n)-w((k+j)h_n)|
\le M(|(k+j+1)h_n-t|+|(k+j)h_n-t|)
\le 8Mh_n.
\]
Zatem $A_n(M)$ zachodzi dla wszystkich dostatecznie
dużych $n$, w sprzeczności z poprzednim wnioskiem.
Przecinając zbiory pełnej miary po $M\in\N$,
wykluczamy wszystkie punkty różniczkowalności wewnątrz
$[0,T]$ naraz.

Dla końca $t=0$, istnienie skończonej pochodnej
prawostronnej oznaczałoby $|W_{h_n}|\le Mh_n$
dla pewnego całkowitego $M$ i wszystkich dużych $n$.
Jednak $\Prob(|W_{h_n}|\le Mh_n)\le C M\sqrt{h_n}$,
co jest sumowalne. Borel--Cantelli wyklucza taką
sytuację. W punkcie $T$ stosujemy analogiczny argument
do $W_T-W_{T-h_n}$. To kończy dowód.
\end{proof}

\begin{wniosek}[Próg wykładnika H\"oldera]
Z prawdopodobieństwem $1$ na żadnym przedziale
$[a,b]\subset[0,\infty)$ o $a<b$ trajektoria Wienera
nie spełnia oszacowania
$|W_t-W_s|\le C|t-s|^\gamma$ dla wszystkich $s,t\in[a,b]$
z jakimkolwiek $C<\infty$ i $\gamma>1/2$.
\end{wniosek}
\begin{proof}
Gdyby takie oszacowanie zachodziło, to dla równomiernego
podziału na $n$ przedziałów suma kwadratów przyrostów
nie przekraczałaby $C^2(b-a)^{2\gamma}n^{1-2\gamma}$,
a więc dążyłaby do zera. To sprzeczne z dodatnią
wariacją kwadratową na każdym przedziale o wymiernych
końcach (przy niemal pewnym wspólnym zdarzeniu).
Każdy przedział $[a,b]$ z $a<b$ zawiera przedział
$[r,s]$ o wymiernych końcach i dodatniej długości;
oszacowanie H\"oldera na $[a,b]$ przeniosłoby się na
$[r,s]$. To daje sprzeczność.
\end{proof}

\begin{center}
\begin{tikzpicture}[x=0.9cm,y=0.8cm,>=Latex]
\draw[->] (-0.2,0)--(10.5,0) node[right] {$t$};
\draw[->] (0,-1.3)--(0,2.6) node[above] {$W_t$};
\draw[blue!75!black,thick] (0,0)--(0.45,0.55)--(0.8,0.18)--(1.18,1.05)--(1.55,0.82)--(1.83,1.5)--(2.14,0.9)--(2.5,1.1)--(2.86,0.22)--(3.2,0.67)--(3.45,0.08)--(3.75,-0.45)--(4.1,0.36)--(4.4,-0.05)--(4.7,0.78)--(5.04,0.42)--(5.37,1.2)--(5.65,0.66)--(5.95,1.3)--(6.25,0.55)--(6.56,0.1)--(6.86,0.88)--(7.21,0.47)--(7.5,1.22)--(7.83,1.0)--(8.15,1.72)--(8.53,1.03)--(8.85,1.4)--(9.2,0.82)--(9.6,1.65)--(10,1.15);
\draw[dashed,red!70!black] (0,1.42)--(10,1.42);
\node[red!70!black,above left] at (10,1.42) {$a$};
\draw[densely dotted,red!70!black] (1.80,0)--(1.80,1.42);
\node[below] at (1.80,0) {$\tau_a$};
\end{tikzpicture}
\par\smallskip\small Schemat trajektorii i pierwszego trafienia w poziom $a$.
Rysunek ma charakter poglądowy: prawdziwa trajektoria Wienera
jest ciągła, ale nie jest łamaną o skończonej liczbie odcinków.
\end{center}

\section{Pierwsze trafienie i zasada odbicia}
\begin{definicja}[Czas pierwszego trafienia]
Niech $W$ będzie ciągłym procesem Wienera i $a>0$.
Definiujemy $\tau_a:=\inf\{t\ge0:W_t=a\}$,
przyjmując $\inf\varnothing=\infty$.
Dla naturalnej filtracji (lub dowolnej filtracji, do której
$W$ jest adaptowany) jest to czas zatrzymania, gdyż
$\{\tau_a\le t\}=\{\sup_{0\le s\le t}W_s\ge a\}$,
a supremum ciągłej trajektorii można obliczać po
chwilach wymiernych z $[0,t]$ wraz z końcem przedziału.
\end{definicja}

\begin{twierdzenie}[Zasada odbicia]
Dla standardowego procesu Wienera $W$, $t>0$ i $a>0$
zachodzi
\[
\Prob\left(\sup_{0\le s\le t}W_s\ge a\right)
=2\Prob(W_t\ge a)
=2\left(1-\Phi\left(\frac{a}{\sqrt t}\right)\right),
\]
gdzie $\Phi$ jest dystrybuantą standardowego rozkładu normalnego.
\end{twierdzenie}
\begin{proof}
Niech $\tau=\tau_a$. Z ciągłości trajektorii
$\{\sup_{s\le t}W_s\ge a\}=\{\tau\le t\}$.
Jeśli $W_t>a$, to z ciągłości i $W_0=0<a$ musieliśmy
wcześniej osiągnąć $a$, zatem $\{W_t>a\}\subset\{\tau\le t\}$.

Na zdarzeniu $\{\tau\le t\}$ mamy $W_\tau=a$.
Z mocnej własności Markowa przyszły przyrost
$W_t-W_\tau$, przy zadanym $\calF_\tau$ i czasie
pozostałym $t-\tau$, ma symetryczny rozkład
$\Normal(0,t-\tau)$. Możemy to sformalizować,
warunkując najpierw względem $\calF_\tau$ i stosując
mocną własność Markowa do przyszłego procesu
$B_u=W_{\tau+u}-W_\tau$ oraz mierzalnej funkcji
$u=t-\tau$; fakt dla prostych $u$ przenosimy do
ogólnego $u$ przez aproksymację i ciągłość trajektorii.
Stąd
\[
\Prob(\tau\le t,\ W_t<a)
=\Prob(\tau\le t,\ W_t>a).
\]
Równość jest prawdziwa także na granicy $\tau=t$,
gdyż oba zdarzenia tam nie zachodzą.
Ponieważ $W_t$ ma rozkład ciągły,
$\Prob(W_t=a)=0$ i
$\Prob(\tau\le t,\ W_t>a)=\Prob(W_t>a)$.
Dzieląc $\{\tau\le t\}$ według zdarzeń
$\{W_t<a\},\{W_t=a\},\{W_t>a\}$,
dostajemy żądany wzór.
\end{proof}

\begin{wniosek}[Rozkład czasu trafienia]
Dla $a>0$ czas $\tau_a$ jest skończony prawie na pewno,
a jego dystrybuanta i gęstość na $(0,\infty)$ mają postać
\[
\Prob(\tau_a\le t)=2\left(1-\Phi(a/\sqrt t)\right),
\qquad
f_{\tau_a}(t)=\frac{a}{\sqrt{2\pi}\,t^{3/2}}
\exp\left(-\frac{a^2}{2t}\right).
\]
Jednak $\E\tau_a=\infty$.
\end{wniosek}
\begin{proof}
Pierwszy wzór jest zasadą odbicia. Gdy $t\to\infty$,
argument $a/\sqrt t$ dąży do zera, więc dystrybuanta
dąży do $2(1-\Phi(0))=1$: czas jest skończony p.n.
Różniczkując dystrybuantę po $t>0$, dostajemy gęstość.
Ponadto, z twierdzenia o wartości średniej dla $\Phi$,
\[
\Prob(\tau_a>t)=2(\Phi(a/\sqrt t)-1/2)
\sim \sqrt{\frac2\pi}\frac{a}{\sqrt t}
\quad\text{gdy }t\to\infty.
\]
Dla nieujemnej zmiennej losowej zachodzi
$\E\tau_a=\int_0^\infty\Prob(\tau_a>t)\,\dd t$
(z twierdzenia Tonellego). Całka jest rozbieżna,
gdyż $\int_1^\infty t^{-1/2}\,\dd t=\infty$.
\end{proof}

\begin{uwaga}[Prawie pewne zakończenie nie daje skończonego średniego czasu]
Ruch Browna dociera do ustalonego dodatniego poziomu
z prawdopodobieństwem $1$, lecz wartość oczekiwana
czasu oczekiwania jest nieskończona.
To ważny kontrprzykład: nie wolno bez dodatkowych
warunków stosować próbkowania opcjonalnego bezpośrednio
w chwili $\tau_a$. Istnienie granicy $W_{\tau_a}=a$
prawie na pewno nie oznacza, że
$\E W_{\tau_a}=\E W_0$.
\end{uwaga}

\begin{zadanie}[Transformata Laplace'a czasu pierwszego trafienia]
Niech $W$ będzie procesem Wienera, $a>0$,
$\tau_a=\inf\{t\ge0:W_t=a\}$.
Obliczyć $\E e^{-\lambda\tau_a}$ dla $\lambda>0$.

\emph{Rozwiązanie.} Dla $\theta=\sqrt{2\lambda}>0$
proces $M_t=\exp(\theta W_t-\theta^2t/2)$
jest martyngałem: wynika to z niezależności przyrostów i
$\E e^{\theta(W_t-W_s)}=e^{\theta^2(t-s)/2}$.
Dla ograniczonego czasu $\tau_a\wedge n$ twierdzenie o
próbkowaniu opcjonalnym daje
$1=\E M_{\tau_a\wedge n}$.
Przed trafieniem w $a$ mamy $W_t<a$, a w chwili
trafienia $W_{\tau_a}=a$, więc
$0\le M_{\tau_a\wedge n}\le e^{\theta a}$.
Wiemy już, że $\tau_a<\infty$ p.n.; zatem
$M_{\tau_a\wedge n}\to\exp(\theta a-\lambda\tau_a)$
p.n. Zbieżność zdominowana pozwala przejść do granicy:
$1=e^{\theta a}\E e^{-\lambda\tau_a}$, skąd
\[
\E e^{-\lambda\tau_a}=e^{-a\sqrt{2\lambda}}.
\]
Ta transformata będzie w przyszłości użyteczna przy
porównywaniu ruchu Browna z procesami Lévy'ego.
\end{zadanie}

\section{Wyjście z przedziału: prawdopodobieństwa i średni czas}
\begin{twierdzenie}[Wyjście procesu Wienera z $(-a,b)$]
Niech $a,b>0$ oraz
$\tau:=\inf\{t\ge0:W_t\notin(-a,b)\}$.
Wtedy $\tau<\infty$ prawie na pewno,
$\E\tau=ab$ oraz
\[
\Prob(W_\tau=b)=\frac{a}{a+b},\qquad
\Prob(W_\tau=-a)=\frac{b}{a+b}.
\]
\end{twierdzenie}
\begin{proof}
Z ciągłości $W_{t\wedge\tau}\in[-a,b]$ i
$W_\tau\in\{-a,b\}$ na zdarzeniu $\{\tau<\infty\}$.
Z przyrostów gaussowskich wynika, że zarówno $W_t$,
jak i $W_t^2-t$ są martyngałami. Dla drugiego procesu
sprawdzamy bowiem, że przy $s<t$:
\[
\E[W_t^2-t\mid\calF_s]
=W_s^2+\E[(W_t-W_s)^2]-t=W_s^2-s,
\]
gdyż mieszany wyraz ma zerową warunkową wartość oczekiwaną.
Z próbkowania opcjonalnego w ograniczonej chwili
$t\wedge\tau$ dostajemy
$\E W_{t\wedge\tau}=0$ oraz
$\E W_{t\wedge\tau}^2=\E(t\wedge\tau)$.
Wprowadzając $M=\max(a,b)$, otrzymujemy
$\E(t\wedge\tau)\le M^2$. Zbieżność monotoniczna
przy $t\to\infty$ daje $\E\tau\le M^2<\infty$,
a więc $\tau<\infty$ p.n.

Ze zbieżności zdominowanej,
$W_{t\wedge\tau}\to W_\tau$ zarówno w $L^1$,
jak i w $L^2$; zatem
$\E W_\tau=0$ i
$\E W_\tau^2=\E\tau$.
Niech $p=\Prob(W_\tau=b)$. Ponieważ
$W_\tau$ przyjmuje tylko dwie wartości,
$0=\E W_\tau=bp-a(1-p)$, skąd $p=a/(a+b)$.
Następnie
\[
\E\tau=\E W_\tau^2
=b^2\frac{a}{a+b}+a^2\frac{b}{a+b}=ab.
\]
Wszystkie przejścia do granicy zostały więc uzasadnione
odpowiednio zbieżnością monotoniczną i zdominowaną.
\end{proof}

\begin{zadanie}[Start z dowolnego punktu przedziału]
Niech $x\in(-a,b)$, $a,b>0$, a proces
$X_t=x+W_t$. Obliczyć prawdopodobieństwo wyjścia
przez prawą krawędź i oczekiwany czas wyjścia.

\emph{Rozwiązanie.} Wyjście $X$ z $(-a,b)$ jest
równoważne wyjściu $W$ z przedziału
$(-a-x,b-x)$. W poprzednim twierdzeniu podstawiamy
$a'=a+x$, $b'=b-x$, obie liczby są dodatnie. Dostajemy
\[
\Prob(X_\tau=b)=\frac{a+x}{a+b},\qquad
\E\tau=(a+x)(b-x).
\]
Zwróćmy uwagę, że funkcja $u(x)=(a+x)(b-x)$
spełnia równanie $\frac12u''(x)=-1$ na $(-a,b)$
i warunki brzegowe $u(-a)=u(b)=0$.
Jest to pierwsza zapowiedź ścisłego związku procesów
Markowa z równaniami różniczkowymi.
\end{zadanie}

\section{Miara Wienera: całkowanie po wszystkich trajektoriach}
Proces $W$ można opisywać dwojako. Dla ustalonej chwili $t$
zmienna $W_t$ jest normalna. Ale pojedynczy wynik
$\omega\in\Omega$ wyznacza także całą ciągłą funkcję
$t\mapsto W_t(\omega)$. Możemy zatem zapytać nie tylko
o prawdopodobieństwo $\{W_t\in A\}$, lecz także o
prawdopodobieństwo zdarzenia zależącego od całej drogi,
na przykład: ``trajektoria przez cały czas pozostaje
poniżej bariery'' lub ``pole pod trajektorią jest dodatnie''.

\begin{definicja}[Przestrzeń trajektorii i miara Wienera]
Dla $T>0$ definiujemy
\[
\calC_0([0,T]):=\{w\in C([0,T];\R):w(0)=0\},
\qquad \norm{w}_\infty=\sup_{0\le t\le T}|w(t)|.
\]
Niech $W$ będzie standardowym procesem Wienera o ciągłych
trajektoriach, traktowanym jako odwzorowanie
$W:\Omega\to\calC_0([0,T])$,
$W(\omega)(t)=W_t(\omega)$.
Jego \emph{prawem na przestrzeni trajektorii}, nazywanym
\emph{miarą Wienera}, jest miara probabilistyczna
$\mu_T:=\Prob\circ W^{-1}$ na
$(\calC_0([0,T]),\calB(\calC_0([0,T])))$.
\end{definicja}

\begin{stwierdzenie}[Mierzalność odwzorowania trajektorii]
Jeżeli $W$ ma ciągłe trajektorie i każda zmienna $W_t$
jest mierzalna, to odwzorowanie
$\omega\mapsto W(\omega)\in\calC_0([0,T])$ jest mierzalne.
\end{stwierdzenie}
\begin{proof}
Dla każdej wymiernej chwili $r\in[0,T]$ odwzorowanie
ewaluacji $e_r(w)=w(r)$ jest ciągłe w normie supremum,
a $e_r\circ W=W_r$ jest mierzalne.
Przestrzeń $\calC_0([0,T])$ jest ośrodkową przestrzenią
metryczną (można przybliżać funkcje ciągłe łamanymi
o wymiernych węzłach i wymiernych wartościach).
Borelowskie sigma-ciało tej przestrzeni jest generowane
przez ewaluacje w przeliczalnym gęstym zbiorze chwil:
rzeczywiście $\norm{w-v}_\infty=\sup_{r\in D}|w(r)-v(r)|$
dla gęstego przeliczalnego $D$ zawierającego końce,
więc każda kula otwarta jest mierzalna względem
$\sigma(e_r:r\in D)$; dzięki ośrodkowości każdy zbiór
otwarty jest przeliczalną sumą takich kul.
Zatem przeciwobrazy zbiorów borelowskich przez $W$
są mierzalne.
\end{proof}

\begin{twierdzenie}[Całki cylindryczne względem miary Wienera]
Niech $T>0$, $0<t_1<\cdots<t_m\le T$ oraz
$F:\R^m\to\R$ będzie funkcją borelowską, nieujemną
lub całkowalną względem rozkładu wektora
$(W_{t_1},\ldots,W_{t_m})$. Wówczas
\begin{align*}
&\int_{\calC_0([0,T])}F(w(t_1),\ldots,w(t_m))\,\dd\mu_T(w)\\
&\quad=\int_{\R^m}F(x_1,\ldots,x_m)
\prod_{k=1}^m p_{t_k-t_{k-1}}(x_k-x_{k-1})
\,\dd\lambda_m(x_1,\ldots,x_m),
\end{align*}
gdzie $t_0=0$, $x_0=0$ i $p_h(z)$ jest gęstością
$\Normal(0,h)$.
\end{twierdzenie}
\begin{proof}
Z definicji miary obrazowej całka po lewej jest równa
$\E F(W_{t_1},\ldots,W_{t_m})$.
Przeprowadzamy liniowe przekształcenie od przyrostów
$Z_k=W_{t_k}-W_{t_{k-1}}$ do kolejnych położeń
$x_k=Z_1+\cdots+Z_k$.
Przyrosty są niezależne, $Z_k$ mają gęstości
$p_{t_k-t_{k-1}}$, więc wektor $Z$ ma gęstość
będącą ich iloczynem. Macierz przekształcenia z $Z$
do $x$ jest trójkątna z jedynkami na przekątnej,
a zatem jej wyznacznik ma moduł $1$.
Wzór na zmianę zmiennych w całce Lebesgue'a
daje tezę. Dla $F\ge0$ dopuszczalność całkowania
wynika z twierdzenia Tonellego; w przypadku funkcji
całkowalnej stosujemy rozkład na części dodatnią i ujemną.
\end{proof}

\begin{zadanie}[Całka z całej trajektorii jako zmienna gaussowska]
Niech $A_T:=\int_0^T W_s\,\dd s$, gdzie całka jest zwykłą
całką Riemanna (równoważnie Lebesgue'a) z ciągłej trajektorii.
Obliczyć rozkład $A_T$ i jego funkcję charakterystyczną.

\emph{Rozwiązanie.} Aproksymujmy całkę sumami Riemanna
$A_T^{(n)}=\sum_{k=0}^{n-1}W_{kT/n}\,(T/n)$.
Każda suma jest gaussowska, ponieważ jest liniową
kombinacją wspólnie gaussowskich zmiennych.
Z ciągłości trajektorii $A_T^{(n)}\to A_T$ p.n.
Co więcej, jeśli $\varphi_n(s)$ jest lewym końcem
przedziału siatki zawierającego $s$, to nierówność
Cauchy'ego--Schwarza i twierdzenie Fubiniego dają
\begin{align*}
\E|A_T^{(n)}-A_T|^2
&\le T\int_0^T\E|W_{\varphi_n(s)}-W_s|^2\,\dd s\\
&\le T^2\frac{T}{n}\longrightarrow0.
\end{align*}
Zatem $A_T^{(n)}$ zbiega w $L^2$ do $A_T$,
a jego średnie i wariancje zbiegają do odpowiednich
wielkości granicy. Granica rozkładów normalnych o
zbieżnych parametrach jest normalna (wynika to wprost
ze zbieżności funkcji charakterystycznych).
Ponieważ $\E A_T=0$ i
\[
\Var(A_T)=\int_0^T\!\int_0^T\Cov(W_s,W_t)
\,\dd s\,\dd t
=\int_0^T\!\int_0^T\min(s,t)\,\dd s\,\dd t
=\frac{T^3}{3},
\]
dostajemy $A_T\sim\Normal(0,T^3/3)$ oraz
\[
\E e^{iuA_T}=\exp\left(-\frac{u^2T^3}{6}\right).
\]
W języku miary Wienera jest to dokładnie obliczenie
całki po trajektoriach
\[
\int_{\calC_0([0,T])}
\exp\left(iu\int_0^T w(s)\,\dd s\right)\,\dd\mu_T(w).
\]
\end{zadanie}

\begin{uwaga}[Trzy całki, których nie wolno utożsamiać]
Wyrażenie $\int_0^T W_s\,\dd s$ jest zwykłą całką
po czasie z jednej losowej, ciągłej funkcji.
Wyrażenie $\int_0^T H_s\,\dd W_s$ będzie całką
stochastyczną Itô, skonstruowaną później z przyrostów
procesu Wienera. Wreszcie
$\int_{\calC_0([0,T])}F(w)\,\dd\mu_T(w)$
całkuje funkcjonał $F$ \emph{po wszystkich możliwych
trajektoriach}. Jest to ścisła matematyczna całka Wienera
po drogach, stanowiąca punkt wyjścia do wzoru
Feynmana--Kaca. Oscylacyjna całka Feynmana w czasie
rzeczywistym wymaga odrębnej analizy i nie jest
zwykłą całką względem miary probabilistycznej.
\end{uwaga}


\begin{stwierdzenie}[Most Browna i miara Wienera mają różne prawa]
Dla $T>0$ niech $\mu_T$ będzie miarą Wienera na
$\calC_0([0,T])$, a $\nu_T$ prawem mostu Browna
$\beta_t=W_t-(t/T)W_T$. Wtedy $\mu_T$ i $\nu_T$
są wzajemnie singularne.
\end{stwierdzenie}
\begin{proof}
Zbiór $A=\{w\in\calC_0([0,T]):w(T)=0\}$ jest domknięty,
więc borelowski. Ponieważ $\beta_T=0$, mamy $\nu_T(A)=1$.
Natomiast $W_T\sim\Normal(0,T)$ ma rozkład bezatomowy,
więc $\mu_T(A)=\Prob(W_T=0)=0$. Jest to definicja
wzajemnej singularności miar.
\end{proof}

\begin{uwaga}[Jaki typ zbieżności wariacji kwadratowej udowodniono?]
Dla \emph{dowolnych deterministycznych} podziałów
$[0,T]$ o średnicy zmierzającej do zera wykazaliśmy
zbieżność sum kwadratów przyrostów w $L^2$, a zatem
według prawdopodobieństwa. Zbieżność prawie na pewno
uzyskaliśmy dla \emph{konkretnego ciągu siatek dyadycznych}.
Z tych dwóch zdań nie wolno wyciągać wniosku, że
zbieżność prawie na pewno zachodzi dla każdej dowolnie
wybranej sekwencji podziałów. Ważne jest również,
czy podziały wolno wybierać zależnie od trajektorii.
\end{uwaga}

\section{Równanie ciepła i generator ruchu Browna}
\begin{twierdzenie}[Generator i równanie ciepła]
Niech $(P_t)_{t\ge0}$ będzie półgrupą zdefiniowaną wzorem
$(P_tf)(x)=\E f(x+W_t)$ dla ograniczonych borelowskich
$f:\R\to\R$. Wówczas:
\begin{enumerate}
\item dla każdej $f\in C_b^3(\R)$, czyli funkcji o
ciągłych i ograniczonych pochodnych do rzędu trzeciego,
zachodzi dla każdego $x\in\R$
\[
\lim_{t\downarrow0}\frac{(P_tf)(x)-f(x)}{t}
=\frac12 f''(x);
\]
\item jeżeli $f\in C_b(\R)$, to funkcja
$u(t,x)=(P_tf)(x)$ jest klasy $C^{1,2}$ na
$(0,\infty)\times\R$, spełnia
$\partial_tu=\frac12\partial_{xx}u$
oraz $u(t,x)\to f(x)$ dla każdego $x$ przy $t\downarrow0$.
\end{enumerate}
\end{twierdzenie}
\begin{proof}
Dla $f\in C_b^3(\R)$ wzór Taylora z resztą daje
\[
f(x+h)=f(x)+hf'(x)+\frac12h^2f''(x)+R(x,h),\qquad
|R(x,h)|\le\frac{\norm{f'''}_\infty}{6}|h|^3.
\]
Podstawmy $h=W_t\sim\Normal(0,t)$ i uśrednijmy.
Mamy $\E W_t=0$, $\E W_t^2=t$ i
$\E|W_t|^3=c_3t^{3/2}$. Zatem
$|(P_tf)(x)-f(x)-tf''(x)/2|
\le \norm{f'''}_\infty c_3t^{3/2}/6$.
Po podzieleniu przez $t$ reszta dąży do zera,
co dowodzi pierwszej części.

Dla $t>0$ możemy zapisać
$u(t,x)=\int_\R f(y)p_t(y-x)\,\dd\lambda_1(y)$.
Na każdym przedziale czasowym $t\in[\delta,R]$
($0<\delta<R$) pochodne jądra gaussowskiego
$\partial_t p_t(z)$, $\partial_zp_t(z)$ i
$\partial_{zz}p_t(z)$ są ograniczone przez całkowalne
funkcje zmiennej $z$, niezależne od $t$ w tym przedziale.
Wynika to z jawnej postaci pochodnych, będących
wielomianami w $z$ pomnożonymi przez
$\exp(-z^2/(2R))$, z odpowiednimi stałymi.
Twierdzenie o różniczkowaniu pod znakiem całki daje
istnienie i ciągłość wskazanych pochodnych $u$.
Bezpośredni rachunek gęstości wykazuje
$\partial_t p_t(z)=\frac12\partial_{zz}p_t(z)$,
a więc po przeniesieniu pochodnych pod całkę
$\partial_tu=\frac12\partial_{xx}u$.

Na koniec $W_t\to0$ według prawdopodobieństwa,
bo $\E W_t^2=t\to0$. Dla ustalonego $x$ i
$\varepsilon>0$ ciągłość $f$ w $x$ daje $\delta>0$
takie, że $|f(x+z)-f(x)|<\varepsilon$ dla $|z|<\delta$.
Dzieląc wartość oczekiwaną według $|W_t|<\delta$
i jego dopełnienia, otrzymujemy
$|P_tf(x)-f(x)|\le\varepsilon
+2\norm f_\infty\Prob(|W_t|\ge\delta)\to\varepsilon$.
Dowolność $\varepsilon$ daje wymaganą zbieżność.
\end{proof}

\begin{przyklad}[Od losowej trajektorii do równania różniczkowego]
Dla $f(x)=\cos(kx)$, $k\in\R$, mamy
$P_tf(x)=\E\cos(kx+kW_t)
=e^{-k^2t/2}\cos(kx)$.
Rzeczywiście, część urojona funkcji charakterystycznej
rozkładu $\Normal(0,t)$ znika, a część rzeczywista
wynosi $e^{-k^2t/2}$. Jednocześnie
$\partial_tu=-(k^2/2)u$ i
$\frac12\partial_{xx}u=-(k^2/2)u$.
Jeden rachunek probabilistyczny dostarcza więc
jawnego rozwiązania równania ciepła.
\end{przyklad}

\section{Podsumowanie i kierunek dalszej teorii}
Proces Wienera łączy zaskakujące fakty: ma ciągłe
trajektorie, ale żadna z nich nie jest różniczkowalna;
ma nieskończoną wariację całkowitą, lecz jego wariacja
kwadratowa jest dokładnie równa upływającemu czasowi.
Jednocześnie można go analizować za pomocą elementarnych
narzędzi gaussowskich, mocnej własności Markowa i
równania ciepła.

Przygotowaliśmy również przestrzeń trajektorii i miarę Wienera.
Dalsze rozdziały rozwiną całki po drogach, ruch Browna
w wielu wymiarach i ścisły rachunek stochastyczny.
W szczególności wzór Feynmana--Kaca połączy
ważone miary trajektorii z półgrupami operatorów
Schrödingera, a całka Itô zbuduje matematyczne
narzędzia do równań różniczkowych stochastycznych.


% ==============================================================
% ROZDZIAŁ 8. ROZSZERZENIE TEORII RUCHU BROWNA
% ==============================================================
